
\begin{longtable}{|p{40mm}|>{\RaggedRight}X|>{\RaggedRight}X|}

% Заголовок таблицы (отображается на первой странице)
\caption{Сравнительный анализ подходов к целевой генерации молекул}
\label{tab:final_comparison} \\

\hline
\textbf{Критерий} & \textbf{Оптимизация в латентном пространстве} & \textbf{Прямая оптимизация политики} \\ 
\hline
\endfirsthead

% Повторяющийся заголовок (для всех страниц, кроме первой)
\multicolumn{3}{l}{\small\textit{Продолжение таблицы \ref{tab:final_comparison}}} \\
\hline
\textbf{Критерий} & \textbf{Оптимизация в латентном пространстве} & \textbf{Прямая оптимизация политики} \\ 
\hline
\endhead

% Текст внизу каждой страницы, кроме последней
\hline
\multicolumn{3}{r}{\small\textit{Продолжение на следующей странице}} \\
\endfoot

% Текст в самом конце таблицы
\hline
\endlastfoot

% ----- Тело таблицы -----


Фундаментальный принцип & Двухэтапный подход: (1) Построение низкоразмерного непрерывного представления (латентного пространства) на основе обучающего набора данных. (2) Последующая оптимизация целевых свойств путем поиска в данном пространстве. & Итеративная оптимизация генеративной модели (политики) с целью максимизации функции вознаграждения. Модель обучается генерировать образцы, соответствующие целевым свойствам, в процессе взаимодействия со средой. \\
\addlinespace
\hline
Разнообразие генерируемых структур & Высокое. Модель обучается аппроксимировать распределение данных из обучающей выборки, что способствует генерации химически разнообразных структур при сэмплировании из латентного пространства. & Потенциально низкое. Склонен к эксплуатации найденных решений и коллапсу мод, когда модель генерирует ограниченный набор схожих структур с высоким вознаграждением. Требует введения специальных механизмов для поддержания разнообразия. \\
\addlinespace
\hline
Требования к данным & Требуется крупный и разнообразный набор данных для обучения репрезентативного латентного пространства. Для этапа целевой оптимизации необходим дополнительный набор данных с аннотированными свойствами для обучения прокси-модели. & Обычно используется предобученная на общем корпусе модель. Для целевой оптимизации не требуется специфический набор данных, так как модель обучается в режиме онлайн, генерируя образцы и получая обратную связь от функции вознаграждения. \\
\addlinespace
\hline
Гибкость (интеграция новых свойств) & Низкая. Интеграция нового целевого свойства для оптимизации, как правило, требует переобучения прокси-модели или модификации и переобучения основной архитектуры (например, с условным латентным пространством). & Высокая. Новые целевые свойства могут быть легко интегрированы путем добавления новых компонентов в составную функцию вознаграждения без необходимости изменения архитектуры генеративной модели. \\
\addlinespace
\hline
Интерпретируемость & Относительно высокая. Структура латентного пространства допускает визуализацию, интерполяцию и анализ, что позволяет исследовать взаимосвязи между структурой молекул и их представлением. & Низкая. Процесс принятия решений генеративной моделью (политикой) трудно интерпретируем. Оптимизация направлена на максимизацию метрики, а не на построение семантически осмысленного представления. \hline


% ----- Конец тела таблицы -----

\end{longtable}